\section{Los resultados deben leerse desglosados por dirección, recursos y métrica}

Una vez entrenado el modelo, lo más riesgoso es reportar solo una puntuación promedio. Los doscientos idiomas incluyen into-English, out-of-English y non-English directions, así como high, low y very-low resource groups; el promedio permite que los grupos más numerosos o con puntuaciones más altas oculten las direcciones realmente débiles. En la clase se usan tres tablas para ejercitar paso a paso esta lectura desglosada.

\subsection{Primero, comparar con los sistemas multilingües de la generación anterior}

Esta sección pasa del mecanismo de entrenamiento a los resultados a nivel de sistema y comienza por establecer una línea base histórica. La tabla comparativa de FLORES-101 coloca a NLLB junto a sistemas anteriores como M2M-100 y DeepNet en el mismo evaluation setup. Al leer la tabla, conviene confirmar si coinciden el número de idiomas que cubre cada modelo, los parámetros y datos de entrenamiento, y la metric; solo después se observa el cambio en el promedio.

\lecturefigure{slide-31-flores101-comparison.jpg}{Comparación de NLLB con sistemas multilingües anteriores en FLORES-101}{Video de la clase de Stanford 00:32:06}

\begin{knowledgebox}{Lectura de la figura: primero las definiciones de columna y los faltantes, después el ganador}
La tabla presenta por separado los resultados into-English, out-of-English, non-English y el promedio; `-/-` indica que algunas fuentes no tienen resultados directamente comparables. La mejora de NLLB respalda la idea de una «mejora conjunta de datos, modelo y entrenamiento», pero una sola tabla no basta para aislar la contribución independiente de cada componente.
\end{knowledgebox}

\begin{warningbox}{La comparación entre generaciones de modelos debe controlar el evaluation protocol}
El tokenizador, el conjunto de idiomas, el checkpoint, la configuración de decodificación y la metric implementation alteran la puntuación. Si los protocolos difieren, los valores de la tabla solo sirven como referencia a nivel de sistema en el momento de la clase, y no deben reordenarse como un leaderboard permanente fuera de su contexto.
\end{warningbox}

\subsection{FLORES-200: desglose por nivel de recursos}

La tabla de FLORES-200 divide además out-of-English e into-English en all, high, low y very-low, y presenta non-English en una columna aparte. Así el lector puede ver si la mejora global llega de verdad a la cola de bajos recursos, en lugar de provenir solo de los idiomas con muchos recursos.

\lecturefigure{slide-32-flores200-results.jpg}{Resultados de FLORES-200 reportados por dirección y nivel de recursos}{Video de la clase de Stanford 00:33:09}

\begin{knowledgebox}{Lectura de la figura: primero comparar horizontalmente los niveles de recursos, luego verticalmente las direcciones}
En una misma fila, la caída de high a very-low muestra el efecto de la disponibilidad de recursos; dentro de un mismo nivel de recursos, la diferencia entre into y out-of-English muestra la asimetría de dirección; el agregado non-English depende de la transfer multilingüe. El promedio resume el sistema, pero no sustituye estos tres diagnósticos.
\end{knowledgebox}

En los resultados de la clase aparecen tanto BLEU como chrF++. BLEU compara principalmente la n-gram precision tokenizada y aplica una penalización por longitud; chrF++ se apoya más en n-gramas de caracteres e incorpora word n-grams, por lo que suele ser más robusta en idiomas morfológicamente ricos o con tokenización inestable. Ambas son solo reference-based automatic metrics.

\begin{knowledgebox}{Términos clave: las métricas de las tablas de resultados}
\begin{tabularx}{\textwidth}{@{}L{0.2\textwidth}L{0.27\textwidth}X@{}}
\toprule
Métrica/grupo & Qué problema resuelve & Qué no puede demostrar \\
\midrule
BLEU / spBLEU & Compara coincidencias de n-gramas con una tokenización unificada & No representa directamente el significado, la naturalidad ni la seguridad \\
chrF++ & Reduce la sensibilidad a la tokenización mediante n-gramas de caracteres y de palabras & Sigue dependiendo de la traducción de referencia y no cubre todas las expresiones válidas \\
Into-English & Entrada en múltiples idiomas hacia el inglés & No representa la calidad de generación cuando el destino es un idioma de bajos recursos \\
Out-of-English & Entrada en inglés hacia múltiples idiomas & Se ve afectada fácilmente por la generación en el lado de destino y la escasez de datos \\
Non-English & Traducción entre dos idiomas distintos del inglés & Suele depender de la transfer; la cobertura de direcciones en el entrenamiento es desigual \\
\bottomrule
\end{tabularx}
\end{knowledgebox}

\subsection{Los límites correctos de la comparación con sistemas comerciales}

En la clase se colocan además las direcciones de bajos recursos junto a sistemas externos como Google Translate. Esta comparación sirve para mostrar la competitividad de NLLB en cobertura de investigación y como modelo público, pero los productos externos se actualizan constantemente y sus idiomas soportados, su decodificación y sus datos internos no son visibles, por lo que no constituye una referencia estable.

\lecturefigure{slide-33-external-system-comparison.jpg}{Comparación de resultados entre NLLB y los sistemas comerciales de traducción en el momento de la clase}{Video de la clase de Stanford 00:33:27}

\begin{knowledgebox}{Lectura de la figura: una comparación con productos es, ante todo, una instantánea con fecha}
Primero hay que confirmar qué direcciones de idioma soportan ambos sistemas y después comparar las columnas de bajos recursos; un hueco no debe interpretarse como una puntuación de cero. Los sistemas comerciales pueden usar datos no publicados y reglas manuales, mientras que NLLB pone el énfasis en activos de investigación abiertos. La tabla respalda que NLLB «alcanza una calidad competitiva en algunas direcciones», no que «supere al producto en todos los escenarios».
\end{knowledgebox}

\begin{importantbox}{El 44\% en BLEU es una mejora relativa}
Si la puntuación anterior es $B_{\text{old}}$ y la nueva es $B_{\text{new}}$, la mejora relativa se define como
\[
r=\frac{B_{\text{new}}-B_{\text{old}}}{B_{\text{old}}}\times100\%.
\]
$r=44\%$ no significa que BLEU haya aumentado 44 puntos absolutos ni que cada par de idiomas haya mejorado 44\%. Es la mejora relativa global que reporta el proyecto NLLB bajo el conjunto de comparación y el protocolo especificados.
\end{importantbox}

\teachervoice{Fan ofrece un criterio práctico de la clase: con algunas métricas, alrededor de 30 el sistema empieza a ser «bastante utilizable». Ella lo llama rule of thumb, no un umbral universal; el idioma, la dirección, el dominio y el costo de los errores modifican la utilidad.}

\subsection{Resumen de la sección}

Los resultados deben desglosarse por translation direction, resource level y metric. Las tablas automáticas permiten comparar sistemas con rapidez, pero no demuestran la calidad percibida por el usuario, la seguridad ni una mejora uniforme en todas las direcciones; el 44\% es una mejora relativa sujeta a condiciones, no una garantía de calidad absoluta.
